\documentclass[aspectratio=169,10pt]{beamer}
% Shared preamble for the Algorithms and Programming Concepts slides (Beamer, Metropolis)
\usetheme[progressbar=frametitle,numbering=fraction,background=light,block=fill,sectionpage=none]{metropolis}
\usepackage[T1]{fontenc}
\usepackage{FiraSans}
\usepackage[scaled=.85]{FiraMono}
\usepackage{amsmath,amssymb}
\usefonttheme[onlymath]{serif}
\usepackage{booktabs,array,tabularx}
\usepackage{tikz,pgfplots}
\usepackage{xcolor}
\usepackage{listings}
\usepackage{adjustbox}
\usepackage{fontawesome5}
\hyphenpenalty=10000\exhyphenpenalty=10000
\newenvironment{sfig}[1][0.80]{\begin{adjustbox}{max width=\linewidth,max totalheight=#1\textheight,center}}{\end{adjustbox}}
\usetikzlibrary{arrows.meta,positioning,shapes.geometric,shapes.misc,fit,calc,backgrounds,decorations.pathreplacing,patterns,matrix,chains}
\pgfplotsset{compat=1.17}

\definecolor{tdInk}{HTML}{1D2330}
\definecolor{tdMuted}{HTML}{5D6678}
\definecolor{tdLine}{HTML}{D9D5CA}
\definecolor{tdPaper}{HTML}{FBFAF7}
\definecolor{tdBlue}{HTML}{2457D6}
\definecolor{tdGreen}{HTML}{17855A}
\definecolor{tdAmber}{HTML}{B7791F}
\definecolor{tdRed}{HTML}{C8362F}
\definecolor{tdPurple}{HTML}{7A4FD1}
\setbeamercolor{normal text}{fg=tdInk,bg=tdPaper}
\setbeamercolor{alerted text}{fg=tdRed}
\setbeamercolor{example text}{fg=tdGreen}
\setbeamercolor{frametitle}{bg=tdInk,fg=white}
\setbeamercolor{progress bar}{fg=tdBlue,bg=tdLine}
\setbeamercolor{title separator}{fg=tdBlue}
\setbeamercolor{block title}{bg=tdBlue!12,fg=tdInk}
\setbeamercolor{block body}{bg=tdBlue!5,fg=tdInk}
\setbeamercolor{block title alerted}{bg=tdRed!12,fg=tdRed!80!black}
\setbeamercolor{block body alerted}{bg=tdRed!5,fg=tdInk}
\setbeamercolor{block title example}{bg=tdGreen!14,fg=tdGreen!70!black}
\setbeamercolor{block body example}{bg=tdGreen!6,fg=tdInk}
\setbeamerfont{frametitle}{size=\large,series=\bfseries}
\makeatletter\setlength{\metropolis@frametitle@padding}{1.15ex}\makeatother
\setbeamertemplate{navigation symbols}{}
\setbeamerfont{block title}{size=\small}

% ---- code listings (Python 3.14) ----
\lstdefinestyle{py}{language=Python,
  basicstyle=\ttfamily\fontsize{8.4}{10.1}\selectfont\color{tdInk},
  keywordstyle=\color{tdBlue}\bfseries, commentstyle=\color{tdMuted}\itshape, stringstyle=\color{tdGreen},
  morekeywords={match,case,type,self,True,False,None},
  emph={print,len,range,sorted,min,max,sum,enumerate,zip,heappush,heappop,cache},emphstyle=\color{tdPurple},
  showstringspaces=false, columns=fullflexible, keepspaces=true, frame=single, rulecolor=\color{tdLine},
  backgroundcolor=\color{white}, xleftmargin=2pt, xrightmargin=2pt, aboveskip=2pt, belowskip=0pt, framesep=3pt}
\lstdefinestyle{out}{basicstyle=\ttfamily\fontsize{8}{9.6}\selectfont\color{tdInk}, columns=fullflexible, keepspaces=true,
  frame=single, rulecolor=\color{tdGreen!60}, backgroundcolor=\color{tdGreen!6}, xleftmargin=2pt, xrightmargin=2pt, aboveskip=2pt, belowskip=0pt, framesep=3pt}
% \pycode[size]{snippet} shows snippets/<snippet>.py ; \pyout[size]{snippet} shows snippets/<snippet>.out
\newcommand{\pycode}[2][\relax]{{#1\lstinputlisting[style=py]{snippets/#2.py}}}
\newcommand{\pyout}[2][\relax]{{#1\lstinputlisting[style=out]{snippets/#2.out}}}

% ---- tikz styles ----
\tikzset{>={Stealth[length=2.2mm]},
  card/.style={rounded corners=4pt,draw=#1!70!black,fill=#1!10,align=center,inner sep=5pt,font=\small}, card/.default=tdBlue,
  pill/.style={rounded corners=7pt,draw=none,fill=#1,text=white,inner sep=3pt,font=\scriptsize\bfseries,align=center}, pill/.default=tdBlue,
  arr/.style={->,thick,draw=tdMuted},
  cell/.style={draw=tdInk!70,minimum width=7.5mm,minimum height=7.5mm,inner sep=0pt,font=\small\ttfamily,fill=white},
  tn/.style={circle,draw=tdInk!80,fill=white,minimum size=7mm,inner sep=0pt,font=\small\bfseries}}
\newcommand{\key}[1]{\textbf{\color{tdBlue}#1}}
\newcommand{\bad}[1]{\textbf{\color{tdRed}#1}}
\newcommand{\good}[1]{\textbf{\color{tdGreen!70!black}#1}}
\newcommand{\apcfooter}[1]{\setbeamertemplate{footline}{\hfill\tiny\color{tdMuted}Algorithms and Programming Concepts \,|\, #1 \,|\, BUV\hspace{1.6em}\insertframenumber/\inserttotalframenumber\hspace{0.8em}\vspace{4pt}}}
\newcommand{\apctitle}[2]{\title{#1}\subtitle{#2}\author{Dr.\ Quang-Vinh Dang\\ {\small Senior Lecturer (Assoc.\ Prof)}}\institute{British University Vietnam}\date{Algorithms and Programming Concepts \,|\, Year 1 \,|\, Python 3.14}}

\apcfooter{Dynamic programming}
\apctitle{Dynamic programming}{Solving a big problem by remembering the answers to small ones}

\tikzset{ft/.style={circle,draw=tdInk!60,fill=white,minimum size=5.4mm,inner sep=0pt,font=\tiny\bfseries},
  level 1/.style={sibling distance=34mm,level distance=7.5mm},level 2/.style={sibling distance=17mm,level distance=7.5mm},
  level 3/.style={sibling distance=9mm,level distance=7.5mm},level 4/.style={sibling distance=7mm,level distance=7.5mm}}
% a table figure written by the Python demo
\newcommand{\figtab}[2][0.5]{\begin{sfig}[#1]\input{snippets/#2.out}\end{sfig}}

\begin{document}
\begin{frame}[plain]\titlepage\end{frame}

\begin{frame}{Don't solve the same problem twice}
\begin{columns}[T]
\begin{column}{0.5\textwidth}
{\small
Imagine adding up a long column of numbers. Someone interrupts you and asks you to add one more number at the bottom. Do you start from the top again?

\vspace{2mm}
Of course not: you \textbf{remember} the total so far. \key{Dynamic programming (DP)} is exactly this habit, applied to harder problems:
\begin{enumerate}\setlength\itemsep{2pt}
\item break the problem into \textbf{smaller sub-problems};
\item solve each sub-problem \textbf{once} and store the answer;
\item build the final answer from the stored ones.
\end{enumerate}}
\end{column}
\begin{column}{0.46\textwidth}
\begin{block}{After this lecture you can}
{\small
\begin{itemize}\setlength\itemsep{2pt}
\item see when a recursive solution repeats work;
\item fix it with \textbf{memoisation} (\texttt{@cache}) or a \textbf{table};
\item follow the DP recipe: state, recurrence, base case, order;
\item solve \textbf{coin change}, \textbf{LCS}, \textbf{edit distance}, \textbf{knapsack} and more;
\item rebuild the \emph{solution}, not only its value;
\item explain why DP turns exponential into polynomial.
\end{itemize}}
\end{block}
\end{column}
\end{columns}
\end{frame}

\begin{frame}{Fibonacci: a clean recursion with a hidden problem}
\begin{columns}[T]
\begin{column}{0.53\textwidth}
\pycode{fib_naive}
\vspace{2mm}
{\small $F(0)=0,\ F(1)=1,\ F(n)=F(n-1)+F(n-2)$. The code is a direct translation of the definition. It is correct, and beautiful. How fast is it?}
\end{column}
\begin{column}{0.43\textwidth}
\pyout{fib_calls}
\vspace{2mm}
{\small
$F(30)$ takes \textbf{2.7 million} calls to compute a number that needs only 30 additions. Adding 10 to $n$ multiplies the work by more than 100. This is exponential growth, roughly $O(1.6^n)$.

\vspace{2mm}
\emph{Where does the waste come from?}}
\end{column}
\end{columns}
\end{frame}

\begin{frame}{The recursion tree of \texttt{fib(5)}: the same work again and again}
\begin{columns}[T]
\begin{column}{0.56\textwidth}
\begin{sfig}[0.72]
\begin{tikzpicture}[ft]
  \node[ft] {5}
    child{node[ft,fill=tdBlue!12]{4}
      child{node[ft,fill=tdAmber!40]{3}
        child{node[ft,fill=tdPurple!25]{2} child{node[ft,fill=tdRed!15]{1}} child{node[ft]{0}}}
        child{node[ft,fill=tdRed!15]{1}}}
      child{node[ft,fill=tdPurple!25]{2} child{node[ft,fill=tdRed!15]{1}} child{node[ft]{0}}}}
    child{node[ft,fill=tdAmber!40]{3}
      child{node[ft,fill=tdPurple!25]{2} child{node[ft,fill=tdRed!15]{1}} child{node[ft]{0}}}
      child{node[ft,fill=tdRed!15]{1}}};
\end{tikzpicture}
\end{sfig}
\end{column}
\begin{column}{0.42\textwidth}
{\small
\textbf{15 calls} for $n=5$. Count the repeats:
\begin{itemize}\setlength\itemsep{2pt}
\item \colorbox{tdAmber!40}{\texttt{fib(3)}} is computed \textbf{2} times,
\item \colorbox{tdPurple!25}{\texttt{fib(2)}} \textbf{3} times,
\item \colorbox{tdRed!15}{\texttt{fib(1)}} \textbf{5} times.
\end{itemize}
Each repeat re-does an entire subtree. For $n=30$ the tree has millions of nodes but only \textbf{31 different} sub-problems.

\vspace{2mm}
\key{Overlapping sub-problems}: the same small problem appears many times. That is the signal that DP will help.}
\end{column}
\end{columns}
\end{frame}

\begin{frame}{Fix 1: memoisation, or ``write it down''}
\begin{columns}[T]
\begin{column}{0.5\textwidth}
\pycode{fib_memo}
\vspace{2mm}
\pyout{memo_demo}
\end{column}
\begin{column}{0.46\textwidth}
{\small
Same recursion, plus one line: \texttt{@cache} (from \texttt{functools}) stores every result. The second time \texttt{fib\_memo(k)} is called with the same \texttt{k}, the answer is looked up instead of recomputed.

\vspace{2mm}
\texttt{cache\_info()} confirms it: only \textbf{31 misses} (the 31 sub-problems $F(0)\ldots F(30)$ really computed) and 28 hits (answers reused).

\vspace{2mm}
This is the \key{top-down} way: start from the big problem, recurse, and remember. Work: $O(n)$.}
\end{column}
\end{columns}
\end{frame}

\begin{frame}{The call tree with a cache: hits cut the tree}
\begin{columns}[T]
\begin{column}{0.5\textwidth}
\begin{sfig}[0.7]
\begin{tikzpicture}[ft,
  level 1/.style={sibling distance=22mm,level distance=9mm},level 2/.style={sibling distance=13mm,level distance=9mm},
  level 3/.style={sibling distance=12mm,level distance=9mm},level 4/.style={sibling distance=10mm,level distance=9mm},
  hit/.style={ft,fill=tdGreen!40,draw=tdGreen!70!black,dashed},miss/.style={ft,fill=tdBlue!12}]
  \node[miss] {5}
    child{node[miss]{4}
      child{node[miss]{3}
        child{node[miss]{2} child{node[miss]{1}} child{node[miss]{0}}}
        child{node[hit]{1}}}
      child{node[hit]{2}}}
    child{node[hit]{3}};
\end{tikzpicture}
\end{sfig}
\vspace{1mm}
{\small \colorbox{tdBlue!12}{blue}: computed (a \emph{miss}). \colorbox{tdGreen!40}{green}: answer found in the cache (a \emph{hit}); nothing below it is explored.}
\end{column}
\begin{column}{0.46\textwidth}
\pyout[\footnotesize]{memo_trace}
\vspace{2mm}
{\small
The real call order of the cached \texttt{fib(5)}: 9 calls instead of 15, and for $n=30$ it is 59 instead of 2.7 million.

\vspace{1mm}
The left branch does all the work. When the right branch asks for \texttt{fib(3)}, the answer is already stored: \textbf{one lookup, no subtree}.}
\end{column}
\end{columns}
\end{frame}

\begin{frame}{Fix 2: tabulation, or ``build from the bottom''}
\begin{columns}[T]
\begin{column}{0.46\textwidth}
\pycode{fib_table}
\vspace{2mm}
\pyout{table_demo}
\vspace{2mm}
{\small No recursion at all: start with the base cases and fill a table from small to large. This is the \key{bottom-up} way.}
\end{column}
\begin{column}{0.5\textwidth}
\pyout{fib_timing}
\vspace{2mm}
{\small
Memoisation and tabulation give the \emph{same} savings: the exponential tree collapses to $n$ steps. The naive version takes well over 100 ms for $n=32$; the table takes about one microsecond.

\vspace{2mm}
\begin{block}{Top-down or bottom-up?}
{\small Top-down is easier to write (add \texttt{@cache}). Bottom-up avoids recursion limits and is slightly faster.}
\end{block}}
\end{column}
\end{columns}
\end{frame}

\begin{frame}{Climbing stairs: each entry depends on two earlier ones}
\figtab[0.22]{fig_stairs}
\vspace{1mm}
\begin{columns}[T]
\begin{column}{0.5\textwidth}
\pycode{stairs}
\vspace{1mm}
\pyout{stairs_demo}
\end{column}
\begin{column}{0.46\textwidth}
{\small
You may climb \textbf{1 or 2 stairs} at a time. To arrive at stair $n$ you came from $n-1$ or from $n-2$, so
\[\mathrm{ways}(n)=\mathrm{ways}(n-1)+\mathrm{ways}(n-2).\]
The arrows show which two cells stair 7 reads. It is Fibonacci again (shifted by one): 21 ways to reach stair 7, and about 20 billion for stair 50, computed in an instant.}
\end{column}
\end{columns}
\end{frame}

\begin{frame}{The DP recipe}
\renewcommand{\arraystretch}{1.35}\small
\begin{tabular}{@{}>{\raggedright\bfseries}p{42mm}>{\raggedright}p{50mm}>{\raggedright\arraybackslash}p{50mm}@{}}
\toprule
Step & Fibonacci & Coin change (next)\\\midrule
1. \textbf{State}: what does one table entry mean? & \texttt{F[i]} = the $i$-th Fibonacci number & \texttt{best[a]} = fewest coins to make amount $a$\\
2. \textbf{Recurrence}: how does a state use smaller ones? & $F[i]=F[i-1]+F[i-2]$ & $\mathrm{best}[a]=1+\min_c \mathrm{best}[a-c]$\\
3. \textbf{Base case}: where does it start? & $F[0]=0,\ F[1]=1$ & $\mathrm{best}[0]=0$\\
4. \textbf{Order}: which entries first? & small $i$ before large $i$ & small amounts first\\
5. \textbf{Answer}: which entry do we want? & \texttt{F[n]} & \texttt{best[amount]}\\
\bottomrule
\end{tabular}
\vspace{4mm}

{\small \textbf{Two conditions} for DP to work: \key{overlapping sub-problems} (we reuse answers) and \key{optimal substructure} (an optimal answer is built from optimal answers to smaller pieces).}
\end{frame}

\begin{frame}{Coin change: greedy looks right and can be wrong}
\begin{columns}[T]
\begin{column}{0.5\textwidth}
{\small
\textbf{Problem.} Given coin values and an amount, use the \emph{fewest coins}.

\vspace{2mm}
\textbf{Greedy idea:} always take the largest coin that fits.
}
\vspace{1mm}
\pycode{greedy}
\end{column}
\begin{column}{0.46\textwidth}
\pyout{coin_greedy}
\vspace{2mm}
{\small
With coins $\{1,3,4\}$ and amount $6$, greedy takes $4+1+1$ (3 coins), but $3+3$ uses only \textbf{2}.

\vspace{2mm}
Greedy fails because its first choice can block a better combination. For ordinary currencies (1, 2, 5, 10, \ldots) greedy happens to work, but you cannot rely on it in general.

\vspace{2mm}
DP never commits too early: it tries every coin for every amount.}
\end{column}
\end{columns}
\end{frame}

\begin{frame}{Coin change with DP: filling the table step by step}
{\small Coins $\{1,2,5\}$, amount $11$. $\mathrm{best}[a]=1+\min_{c\le a}\mathrm{best}[a-c]$, $\ \mathrm{best}[0]=0$.}

\vspace{1mm}
{\footnotesize\textbf{Step 1.} The base case and the first small amounts are filled.}\\
\figtab[0.13]{fig_coin1}

\vspace{1mm}
{\footnotesize\textbf{Step 2.} For $a=6$ look back by each coin: $a-1=5$, $a-2=4$, $a-5=1$. Best: $1+\mathrm{best}[5]=2$ or $1+\mathrm{best}[1]=2$ (green).}\\
\figtab[0.17]{fig_coin2}

\vspace{1mm}
{\footnotesize\textbf{Step 3.} The finished table. $\mathrm{best}[11]=1+\min(\mathrm{best}[10],\mathrm{best}[9],\mathrm{best}[6])=1+\min(2,3,2)=3$.}\\
\figtab[0.17]{fig_coin3}
\end{frame}

\begin{frame}[allowframebreaks]{Coin change in Python}
\pycode{coins}
\vspace{2mm}
{\small
\texttt{best[a]} is the table; \texttt{choice[a]} remembers \emph{which coin} gave the best answer. Each of the $A$ amounts tries $k$ coins: $O(A\cdot k)$. Brute force would try all combinations: exponential. \texttt{INF} marks amounts that cannot be made.}
\framebreak
\pycode{coin_path}
\vspace{2mm}
\pyout{coin_demo}
\vspace{2mm}
{\small
To list the coins, start at the target amount and repeatedly subtract the stored choice: $11\to6\to1\to0$ uses $5,5,1$.

\vspace{1mm}
\textbf{Tip:} always store a \emph{choice} (or parent) next to the value, so you can rebuild the solution, not just its size.}
\end{frame}

\begin{frame}{Longest common subsequence (LCS)}
\begin{columns}[T]
\begin{column}{0.5\textwidth}
{\small
A \textbf{subsequence} keeps the order but may skip letters. The LCS of two strings is the longest sequence that is a subsequence of both.

\vspace{2mm}
\texttt{A\,\textcolor{tdBlue}{B}\,\textcolor{tdBlue}{C}\,\textcolor{tdBlue}{B}\,D\,\textcolor{tdBlue}{A}\,B}\ and \texttt{\textcolor{tdBlue}{B}\,D\,\textcolor{tdBlue}{C}\,A\,\textcolor{tdBlue}{B}\,\textcolor{tdBlue}{A}}\ share \textbf{BCBA}, length 4.

\vspace{2mm}
\textbf{Where is it used?} \texttt{diff} and \texttt{git diff} (which lines did not change?), comparing DNA sequences, plagiarism checks, spell correction.}
\end{column}
\begin{column}{0.46\textwidth}
\begin{block}{The recurrence}
{\small Let $t[i][j]$ = LCS length of the first $i$ letters of $a$ and the first $j$ letters of $b$.
{\footnotesize\[
t[i][j]=\begin{cases}
t[i-1][j-1]+1 & a_i=b_j\\[2pt]
\max(t[i-1][j],\,t[i][j-1]) & \text{else}
\end{cases}\]}
with $t[0][j]=t[i][0]=0$.}
\end{block}
\vspace{1mm}
{\small Letters match: extend the best answer for both shorter strings. They differ: drop the last letter of one string or the other and keep the better.}
\end{column}
\end{columns}
\end{frame}

\begin{frame}{The LCS grid: filling it row by row}
\begin{columns}[T]
\begin{column}{0.52\textwidth}
\figtab[0.78]{fig_lcs1}
\end{column}
\begin{column}{0.44\textwidth}
{\small
The table is filled row by row. Here we stand at the amber cell: $a_3=\texttt C$ and $b_3=\texttt C$ \textbf{match}, so the cell is the \colorbox{tdGreen!35}{diagonal} neighbour plus 1: $t[3][3]=t[2][2]+1=2$.

\vspace{2mm}
If the letters had differed, the cell would be the larger of the two \colorbox{tdBlue!14}{blue} neighbours (above, left).

\vspace{2mm}
Each cell needs just three neighbours, so the whole table costs $O(n\cdot m)$ steps for strings of length $n$ and $m$. Empty cells are not computed yet.}
\end{column}
\end{columns}
\end{frame}

\begin{frame}{The LCS table, and how to read the answer back}
\begin{columns}[T]
\begin{column}{0.52\textwidth}
\figtab[0.78]{fig_lcs2}
\end{column}
\begin{column}{0.44\textwidth}
{\small
\textbf{Traceback} (amber path, from the bottom-right corner): at a \colorbox{tdGreen!40}{green} cell the letters matched, so that letter is in the LCS and we move diagonally; otherwise we move up or left towards the larger neighbour.

\vspace{2mm}
Reading the green cells bottom to top gives \textbf{B, C, B, A}.

\vspace{2mm}
The bottom-right entry, 4, is the length. The path is the \emph{solution}. The figure is drawn from the real table printed by the code on the next slide.}
\end{column}
\end{columns}
\end{frame}

\begin{frame}[allowframebreaks]{LCS in Python}
\pycode{lcs}
\vspace{2mm}
{\small
\texttt{lcs\_table} fills the grid with two loops. A match adds 1 to the diagonal; otherwise the cell takes the larger of the cell above and the cell to the left. Time and memory $O(n\cdot m)$.}
\framebreak
\begin{columns}[T]
\begin{column}{0.54\textwidth}
\pycode{lcs_trace}
\end{column}
\begin{column}{0.42\textwidth}
\pyout{lcs_demo}
\end{column}
\end{columns}
\vspace{1mm}
{\small
\texttt{lcs\_string} starts at the corner and walks back as in the figure, collecting the matched letters in reverse. The printed table is the one drawn two slides earlier.}
\end{frame}

\begin{frame}{Edit distance: how far apart are two words?}
\begin{columns}[T]
\begin{column}{0.5\textwidth}
{\small
The \key{edit distance} (Levenshtein) is the fewest single-letter \textbf{insertions, deletions and replacements} that turn one word into another. Spell checkers, search suggestions (``did you mean\ldots'') and DNA alignment use it.

\vspace{2mm}
\texttt{kitten} $\to$ \texttt{sitting}:
\begin{enumerate}\setlength\itemsep{1pt}
\item \texttt{\textbf{k}itten} $\to$ \texttt{\textbf{s}itten} (replace k by s)
\item \texttt{sitt\textbf{e}n} $\to$ \texttt{sitt\textbf{i}n} (replace e by i)
\item \texttt{sittin} $\to$ \texttt{sittin\textbf{g}} (insert g)
\end{enumerate}
Three edits; no way with two.}
\end{column}
\begin{column}{0.46\textwidth}
\begin{block}{The recurrence}
{\small $d[i][j]$ = distance between the first $i$ letters of $a$ and the first $j$ of $b$.
{\footnotesize\[
d[i][j]=\min\begin{cases}
d[i-1][j]+1 & \text{del}\\
d[i][j-1]+1 & \text{ins}\\
d[i-1][j-1]+[a_i\ne b_j] & \text{sub}
\end{cases}\]}
Base: $d[i][0]=i$, $d[0][j]=j$. \emph{del} = delete $a_i$, \emph{ins} = insert $b_j$.}
\end{block}
\vspace{1mm}
{\small Same shape as LCS: a 2-D table and three neighbours.}
\end{column}
\end{columns}
\end{frame}

\begin{frame}{The edit-distance grid and its traceback}
\begin{columns}[T]
\begin{column}{0.56\textwidth}
\figtab[0.8]{fig_edit}
\end{column}
\begin{column}{0.4\textwidth}
{\small
Row 0 and column 0 are the base cases: turning a word into the empty word costs one deletion per letter.

\vspace{2mm}
The \colorbox{tdAmber!40}{amber path} starts at the bottom-right cell, the answer $d=3$, and walks back to the corner. A \emph{diagonal} step is a keep or a replace, a step \emph{up} is a deletion, a step \emph{left} is an insertion. Where the value rises along a diagonal step, that is a replacement.

\vspace{2mm}
The path is what the code prints on the next slides: \texttt{k->s,e->i,ins g}.}
\end{column}
\end{columns}
\end{frame}

\begin{frame}{Edit distance in Python: the table}
\begin{columns}[T]
\begin{column}{0.54\textwidth}
\pycode{edit}
\end{column}
\begin{column}{0.42\textwidth}
\pyout{edit_demo}
\vspace{2mm}
{\small
The first two loops write the base cases, the double loop applies the recurrence. Time and memory $O(n\cdot m)$.

\vspace{2mm}
\texttt{cost} is 0 when the letters agree and 1 for a replacement. The demo shows two classic pairs and the case of an empty word (three insertions).}
\end{column}
\end{columns}
\end{frame}

\begin{frame}{Edit distance in Python: the traceback}
\begin{columns}[T]
\begin{column}{0.56\textwidth}
\pycode[\fontsize{7.8}{9.2}\selectfont]{edit_ops}
\end{column}
\begin{column}{0.4\textwidth}
{\small
\texttt{edit\_ops} walks back from the corner. At each cell it picks the move that \emph{explains} the cell's value: a diagonal step (keep, or replace when the letters differ), a step up (delete) or a step left (insert).

\vspace{2mm}
It returns the list of edits, in order, and the path of cells (used for the picture two slides earlier).}
\end{column}
\end{columns}
\end{frame}

\begin{frame}{Counting grid paths with obstacles}
\begin{columns}[T]
\begin{column}{0.46\textwidth}
\figtab[0.6]{fig_grid}
\vspace{1mm}
{\small
A robot walks from the top-left to the bottom-right corner, moving only \textbf{right} or \textbf{down}. Dark cells are blocked. How many routes?}
\end{column}
\begin{column}{0.5\textwidth}
{\small
$\mathrm{ways}[r][c]$ = number of routes to cell $(r,c)$. A route arrives from above or from the left:
\[\mathrm{ways}[r][c]=\mathrm{ways}[r-1][c]+\mathrm{ways}[r][c-1],\]
and a blocked cell has \textbf{0} ways.

\vspace{2mm}
Read the picture: the cell right of the wall has $1+1=2$, the cell below the second wall has $1+2=3$. The last cell collects $3+2=5$.

\vspace{2mm}
Without walls a $4\times5$ grid has $\binom{7}{3}=35$ routes (choose which 3 of the 7 moves go down). Walls remove 30 of them.}
\end{column}
\end{columns}
\end{frame}

\begin{frame}[allowframebreaks]{Grid paths in Python}
\pycode{grid}
\framebreak
\pyout{grid_demo}
\vspace{2mm}
{\small
Rows and columns are processed in order, so \texttt{up} and \texttt{left} are already known. The guards \texttt{if r} and \texttt{if c} avoid reading outside the table. Time and memory $O(\text{rows}\times\text{cols})$.

\vspace{2mm}
\textbf{Variation:} keep only one row at a time to use $O(\text{cols})$ memory. A related puzzle, the \emph{minimum-cost path}, replaces ``add'' by ``min'' plus the cell's cost.}
\end{frame}

\begin{frame}{0/1 knapsack: take it or leave it}
\begin{columns}[T]
\begin{column}{0.5\textwidth}
{\small
You met knapsack in the optimization lecture. Here is the DP \textbf{table} that also tells us \emph{which} items to take.

\vspace{2mm}
Items (weight, value): $(1,1)$, $(3,4)$, $(4,5)$, $(5,7)$. Capacity $7$.

\vspace{2mm}
\textbf{State:} $t[i][c]$ = best value using the first $i$ items with capacity $c$.}
\end{column}
\begin{column}{0.46\textwidth}
\begin{block}{The recurrence}
{\small For item $i$ with weight $w$ and value $v$:
\[
t[i][c]=\max\begin{cases}
t[i-1][c] & \text{skip}\\
t[i-1][c-w]+v & \text{take}
\end{cases}\]
$t[0][c]=0$; \emph{take} only if $w\le c$. Answer: $t[n][W]$.}
\end{block}
\vspace{1mm}
{\small Each cell is one decision: \emph{is item $i$ worth the room it needs?}}
\end{column}
\end{columns}
\end{frame}

\begin{frame}{The knapsack table and the chosen items}
\begin{columns}[T]
\begin{column}{0.56\textwidth}
\figtab[0.82]{fig_knap}
\end{column}
\begin{column}{0.4\textwidth}
{\small
Rows are items (in order), columns are capacities $0\ldots7$. The bottom-right cell is the answer: \textbf{9}.

\vspace{2mm}
\textbf{Reconstruction} (amber path): start at the corner. If the value is the \emph{same} as the cell above, the item was skipped, so move up. If it \emph{differs}, the item was taken (\colorbox{tdGreen!40}{green}): move up and left by its weight.

\vspace{2mm}
Items 3 and 2 are taken: weights $4+3=7$, values $5+4=9$.}
\end{column}
\end{columns}
\end{frame}

\begin{frame}[allowframebreaks]{Knapsack with item recovery in Python}
\pycode{knap_table}
\vspace{2mm}
{\small
\texttt{t[i][c]} starts as ``skip item $i$'' (copy the cell above); if the item fits, we take the better of that and ``take it''. Time and memory $O(n\cdot W)$. The 1-D version of the optimization lecture saves memory but loses the information needed to recover the items. The printed matrix (next page) is the one drawn on the previous slide; \texttt{knap\_pick} walks back from the corner: a changed value means the item was taken.}
\framebreak
\pycode{knap_pick}
\vspace{2mm}
\pyout{knap01_demo}
\end{frame}

\begin{frame}{House robber: never rob two neighbours}
\begin{columns}[T]
\begin{column}{0.5\textwidth}
\figtab[0.3]{fig_rob}
\vspace{1mm}
{\small
Houses in a row hold money. Alarms ring if you rob two \textbf{adjacent} houses. Maximise the loot.

\vspace{1mm}
\textbf{State:} $\mathrm{best}[i]$ = most money from the first $i$ houses.

\vspace{1mm}
$\mathrm{best}[i]=\max(\mathrm{best}[i-1],\ \mathrm{best}[i-2]+m_i)$: skip house $i$, or rob it and the best of the first $i-2$.

\vspace{1mm}
Answer: 19, by robbing houses 0, 2 and 5 (money 2, 9, 8).}
\end{column}
\begin{column}{0.46\textwidth}
\pycode{rob}
\vspace{2mm}
{\small
Only the last two values matter, so two variables are enough: memory $O(1)$. \texttt{take} is the best loot if we rob the current house, \texttt{skip} if we do not.

\vspace{1mm}
To recover \emph{which} houses, we keep the whole table (next slide).}
\end{column}
\end{columns}
\end{frame}

\begin{frame}{House robber: recovering the houses}
\begin{columns}[T]
\begin{column}{0.52\textwidth}
\pycode[\fontsize{7.8}{9.2}\selectfont]{rob_plan}
\end{column}
\begin{column}{0.44\textwidth}
\pyout{rob_demo}
\vspace{2mm}
{\small
\texttt{best[i]} is the table of the previous slide. \textbf{Traceback:} if \texttt{best[i]} equals \texttt{best[i-1]}, house $i-1$ was skipped; otherwise it was robbed and we jump two houses back.}
\end{column}
\end{columns}
\end{frame}

\begin{frame}{Longest increasing subsequence (LIS)}
\begin{columns}[T]
\begin{column}{0.54\textwidth}
\figtab[0.3]{fig_lis}
\vspace{1mm}
{\small
Find the longest subsequence of a list whose values strictly \textbf{increase}. Here $[10,9,2,5,3,7,101,18]$ has the LIS $2,5,7,101$, length 4.

\vspace{1mm}
\textbf{State:} $\mathrm{best}[i]$ = length of the longest increasing subsequence that \emph{ends} at position $i$.

\vspace{1mm}
$\mathrm{best}[i]=1+\max\{\mathrm{best}[j]: j<i,\ a_j<a_i\}$ (or 1 if no such $j$). The answer is the \emph{maximum} over all $i$, not the last entry.}
\end{column}
\begin{column}{0.42\textwidth}
\pyout[\footnotesize]{lis_tails}
\vspace{1mm}
{\footnotesize
\textbf{Faster ($O(n\log n)$).} \texttt{tails[k]} = smallest value that ends an increasing subsequence of length $k+1$. \texttt{bisect\_left} places each $x$: it extends the list or replaces a larger tail. The list stays sorted; its length is the LIS length. (\texttt{tails} is \emph{not} the subsequence.)}
\end{column}
\end{columns}
\end{frame}

\begin{frame}{LIS in Python: the $O(n^2)$ version}
\pycode[\fontsize{7.8}{9.2}\selectfont]{lis_n2}
\vspace{1mm}
{\small
For each $i$ look at every earlier $j$. The test \texttt{best[j] >= best[i]} means ``extending $j$ beats the best chain found so far for $i$''. The chain is rebuilt through the \texttt{prev} links.}
\end{frame}

\begin{frame}{LIS in Python: the $O(n\log n)$ version}
\begin{columns}[T]
\begin{column}{0.5\textwidth}
\pycode{lis_fast}
\end{column}
\begin{column}{0.46\textwidth}
\pyout{lis_demo}
\vspace{2mm}
{\small
\texttt{lis\_length} is $O(n\log n)$ and returns only the length; \texttt{bisect\_left} is a binary search on the sorted list \texttt{tails}. Both versions agree: 4.

\vspace{2mm}
To recover the subsequence itself in $O(n\log n)$ you also store, for each tail, the index of its predecessor (a challenge for you).}
\end{column}
\end{columns}
\end{frame}

\begin{frame}{Partition into two equal sums}
\begin{columns}[T]
\begin{column}{0.56\textwidth}
\figtab[0.6]{fig_part}
\vspace{1mm}
{\small Numbers $[1,5,11,5]$ add up to 22, so we need a subset with sum \textbf{11}. Green = this sum can be made with the numbers so far.}
\end{column}
\begin{column}{0.42\textwidth}
{\small
\textbf{Question:} can a list be split into two groups with the same sum?

\vspace{1mm}
Half of the total is the target $T$. If some subset sums to $T$, the rest does too. So: \emph{is $T$ reachable as a subset sum?}

\vspace{1mm}
\textbf{State:} the set of reachable sums. Adding a number $x$: every old sum $s$ gives a new sum $s+x$. The last row shows $11$ reachable: $1+5+5$.

\vspace{1mm}
Python trick: store the set as the bits of one integer; \texttt{reach << x} shifts every sum by $x$ in a single operation.}
\end{column}
\end{columns}
\end{frame}

\begin{frame}{Partition: the code and a few tests}
\begin{columns}[T]
\begin{column}{0.54\textwidth}
\pycode{partition}
\vspace{2mm}
\pyout{part_demo}
\end{column}
\begin{column}{0.42\textwidth}
{\small
\texttt{reach} is a \textbf{bitset}: bit $s$ is 1 when sum $s$ can be made. Start with only sum 0. For each number, \texttt{reach |= reach << x} adds the new sums.

\vspace{2mm}
An odd total can never split. $[1,2,3,5]$ adds to 11: odd, so \texttt{False}. $[3,3,3,4,5]$ sums to 18, and $4+5=3+3+3=9$.

\vspace{2mm}
Time: $O(n\cdot T)$ bit-steps, and Python's big integers do 64 of them per machine word. This is a \emph{subset-sum} problem, a cousin of knapsack.}
\end{column}
\end{columns}
\end{frame}

\begin{frame}{Word break: can the text be cut into dictionary words?}
\begin{columns}[T]
\begin{column}{0.5\textwidth}
\figtab[0.22]{fig_wb}
\vspace{1mm}
{\small
Dictionary: \texttt{cat, cats, and, sand, dog}. Text: \texttt{catsanddog}. A cell \texttt{T} at position $i$ means \textbf{the first $i$ letters can be cut into words}.

\vspace{1mm}
Position $i$ is \texttt{T} if for some word $w$ the prefix of length $i-|w|$ is \texttt{T} and the text ends there with $w$. Here: 3 (\texttt{cat}), 4 (\texttt{cats}), 7 (\texttt{cat+sand} or \texttt{cats+and}), 10 (\ldots\texttt{dog}). Position 10 is \texttt{T}: success.}
\end{column}
\begin{column}{0.46\textwidth}
\pycode[\fontsize{7.6}{9}\selectfont]{word_break}
\vspace{1mm}
\pyout[\fontsize{7.6}{9}\selectfont]{wb_demo}
\end{column}
\end{columns}
\end{frame}

\begin{frame}{When to use dynamic programming}
\renewcommand{\arraystretch}{1.1}\footnotesize
\begin{tabular}{@{}>{\raggedright\bfseries}p{40mm}>{\raggedright}p{52mm}>{\raggedright\arraybackslash}p{44mm}@{}}
\toprule
Problem & State & Cost\\\midrule
Fibonacci, climbing stairs & \texttt{F[i]} & $O(n)$ instead of exponential\\
Coin change & \texttt{best[amount]} & $O(\text{amount}\times\text{coins})$\\
Longest common subsequence & \texttt{t[i][j]} (prefix of $a$, prefix of $b$) & $O(n\cdot m)$\\
Edit distance & \texttt{d[i][j]} (prefixes) & $O(n\cdot m)$\\
0/1 knapsack & \texttt{t[item][capacity]} & $O(n\cdot W)$\\
Longest increasing subsequence & \texttt{best[i]} (ends at $i$) & $O(n^2)$ or $O(n\log n)$\\
House robber & \texttt{best[i]} (first $i$ houses) & $O(n)$\\
Grid paths & \texttt{ways[r][c]} & $O(\text{rows}\times\text{cols})$\\
Partition, word break & reachable sums, good prefixes & $O(n\cdot T)$, $O(n\cdot\text{words})$\\
Shortest paths, negative edges (Bellman--Ford) & \texttt{dist[v][steps]} & $O(V\cdot E)$\\
\bottomrule
\end{tabular}
\vspace{3mm}

{\small \textbf{Not a fit:} problems with no repeated sub-problems (then plain recursion is already fine), or where the state is too large to store. DP trades \textbf{memory for time}.}
\end{frame}

\begin{frame}{Recap}
\begin{columns}[T]
\begin{column}{0.48\textwidth}
\begin{block}{Take-aways}
{\small
\begin{enumerate}\setlength\itemsep{3pt}
\item DP = \key{solve each sub-problem once} and reuse the answer.
\item Look for \key{overlapping sub-problems} in a recursion tree.
\item \key{Top-down}: recursion + \texttt{@cache}. \key{Bottom-up}: fill a table in the right order.
\item Recipe: \textbf{state}, \textbf{recurrence}, \textbf{base case}, \textbf{order}, \textbf{answer}.
\item Greedy can fail where DP succeeds (coin change with $\{1,3,4\}$).
\item Store the \textbf{choice} too, to rebuild the solution by \emph{traceback}.
\end{enumerate}}
\end{block}
\end{column}
\begin{column}{0.48\textwidth}
\begin{block}{Common table shapes}
{\small
\begin{itemize}\setlength\itemsep{3pt}
\item \textbf{1-D, depends on a few earlier cells}: Fibonacci, stairs, robber, coins.
\item \textbf{1-D, depends on all earlier cells}: LIS, word break.
\item \textbf{2-D, two prefixes}: LCS, edit distance.
\item \textbf{2-D, items $\times$ capacity}: knapsack, partition.
\item \textbf{2-D, positions}: grid paths.
\end{itemize}}
\end{block}
\end{column}
\end{columns}
\end{frame}

\begin{frame}{Challenge problems (no solutions: try them!)}
\begin{columns}[T]
\begin{column}{0.5\textwidth}
{\footnotesize
\begin{enumerate}\setlength\itemsep{3pt}
\item \textbf{Count the ways} to make an amount with given coins (order does not matter). Change \texttt{min\_coins} to add instead of minimise; watch the loop order.
\item \textbf{Minimum path sum}: in a grid of costs, find the cheapest right/down route from corner to corner, and print the route.
\item \textbf{Rod cutting}: a rod of length $n$ can be cut in integer pieces with a price per length. Maximise the total price and list the cuts.
\item \textbf{Longest palindromic subsequence} of a string. Hint: it is the LCS of the string and its reverse.
\item \textbf{House robber II}: the houses stand in a \emph{circle}, so the first and the last are neighbours.
\end{enumerate}}
\end{column}
\begin{column}{0.46\textwidth}
{\footnotesize
\begin{enumerate}\setlength\itemsep{3pt}\setcounter{enumi}{5}
\item \textbf{Unbounded knapsack}: every item may be used any number of times. Which loop in \texttt{knapsack\_dp} changes?
\item \textbf{Matrix-chain multiplication}: choose the order of multiplying $n$ matrices to minimise the scalar products. State: $d[i][j]$ for the sub-chain $i\ldots j$.
\item \textbf{Edit distance with costs}: replace costs 2, insert and delete cost 1. Relate the answer to the LCS length.
\item \textbf{Longest common substring} (contiguous!): what changes in the recurrence?
\item \textbf{Count the LIS}: how many different longest increasing subsequences does $[1,3,5,4,7]$ have?
\end{enumerate}}
\end{column}
\end{columns}
\end{frame}
\end{document}
